\begin{table}[t]
\caption*{
{\fontsize{20}{25}\selectfont  Typology Robustness — H1 — Effect of the invasion on sentiment\fontsize{12}{15}\selectfont } \\ 
{\fontsize{14}{17}\selectfont  Same specification under each time-assignment variant\fontsize{12}{15}\selectfont }
} 
\fontsize{12.0pt}{14.0pt}\selectfont
\begin{tabular*}{\linewidth}{@{\extracolsep{\fill}}llll}
\toprule
 & Probabilistic & Day & Drop \\ 
\midrule\addlinespace[2.5pt]
\multicolumn{4}{l}{Main (Table 1, Model 5)} \\[2.5pt] 
\midrule\addlinespace[2.5pt]
Estimate & -0.217*** & -0.267*** & — \\ 
Std. Error & (0.036) & (0.034) & — \\ 
Bootstrap p & — & — & — \\ 
N & 45,366 & 68,973 & — \\ 
Clusters & — & — & — \\ 
\bottomrule
\end{tabular*}
\begin{minipage}{\linewidth}
\vspace{.05em}
Each column re-estimates the same specification under a different
time-assignment typology: probabilistic imputation of missing publication
times (the primary analysis), the calendar day as the running variable, or
articles without a recorded time dropped. Sample sizes differ across columns
by construction, since the typologies retain different articles and select
their own MSE-optimal bandwidths. The coefficient reported is Invasion. Table 1
Model 5 is an rdrobust fit reported with cluster-robust standard errors, so no
bootstrap p-value applies. Stars follow the bootstrap p-value where one exists.\\
\end{minipage}
\end{table}

